\section{Where the pattern and the condition come from}
\label{sec:ddl}

\subsection{The question the comprehension cannot answer}

The comprehension of \S\ref{sec:cfl} presupposes two things it does not supply.
Its pattern \texttt{@spec} must have a grammar of terms to match against. Its condition \texttt{Chi(spec)} is a claim about
structure, so there must be a notion of what structure a specimen has.

In SQL both come from the schema. \texttt{spec} is a row of a known table with
known column types, and \texttt{chi} is a predicate over those columns. The
schema is the data definition language, and the query language leans on it
entirely.

\flang{}'s schema is the \texttt{Theory} declaration, and the reason to keep
reading is that it is a great deal more than a schema.

\subsection{A \texttt{Theory} declaration, whole}

Here is one, complete~\cite{modulesyntax}. It defines the \rhoc{}
calculus~\cite{MR05} itself --- which is a
useful thing to see early, because it shows that the calculus \flang{} is built
on is not privileged: it is one language definition among others, written in
the same form a developer would use for an invoice format.

\begin{lstlisting}[language=mettail]
Theory RhoCalc() {
    Types {
        Proc;
        Name;
    }
    Terms {
        PZero  . |- "0" : Proc;
        PDrop  . n:Name |- "*" "(" n ")" : Proc;
        POutput. n:Name, q:Proc |- n "!" "(" q ")" : Proc;
        PInput . n:Name, ^x.p:[Name -> Proc] |- n "?" x "." "{" p "}" : Proc;
        PPar   . ps:HashBag(Proc) |- "{" ps.*sep("|") "}" : Proc;
        NQuote . p:Proc |- "@" "(" p ")" : Name;
    }
    Equations {
        (NQuote (PDrop N)) == N;
    }
    Rewrites {
        RComm : (PPar {(PInput n ^x.p), (POutput n q), ...rest})
                  ~> (PPar {(subst ^x.p (NQuote q)), ...rest});
        RDrop : (PDrop (NQuote P)) ~> P;
        RPar  : if S ~> T then (PPar {S, ...rest}) ~> (PPar {T, ...rest});
    }
}
\end{lstlisting}

Four blocks. \texttt{Types} names the syntactic categories. \texttt{Terms} gives
one constructor per line: a label, then the arguments it takes with their sorts,
then --- after the turnstile --- the concrete syntax written in terms of those
argument names, then the category it produces. So the grammar and the
pretty-printer come from the same declaration. \texttt{Equations} says which
distinct terms are to be regarded as the same. \texttt{Rewrites} says how terms
move --- in the style Plotkin introduced for operational
semantics~\cite{Plotkin81} and Milner used to such effect for
processes~\cite{Milner92}.

From that block, the toolchain generates the abstract syntax types, a parser and
a bidirectional printer, capture-avoiding substitution that respects binding,
and a rewrite engine.

Three details in the \texttt{Terms} block repay a second look, because each is
carrying weight a conventional schema language has no way to carry.
\texttt{\^{}x.p:[Name -> Proc]} declares that \texttt{PInput}'s second argument
is an \emph{abstraction}: \texttt{x} binds in \texttt{p}. Binding is part of the
data declaration, which is why substitution can be generated rather than
hand-written. \texttt{ps:HashBag(Proc)} declares that \texttt{PPar} takes an
unordered collection with multiplicity, and \texttt{ps.*sep("|")} says how to
render one. And in \texttt{RComm}, \texttt{...rest} binds whatever else was in
the bag and puts it back --- which is what lets a rule fire on two participants
inside an arbitrary parallel composition without saying anything about the
others.

\subsection{The four rungs}
\label{sec:rungs}

The blocks form a ladder, and where a given language definition stops on
that ladder determines what kind of thing it is.

\begin{center}
\small
\begin{tabular}{@{}L{0.26\textwidth}L{0.32\textwidth}L{0.32\textwidth}@{}}
\toprule
\textbf{Blocks populated} & \textbf{What you can present} & \textbf{The familiar name} \\
\midrule
\texttt{Types}, \texttt{Terms} &
free term algebras &
algebraic data types \\[3pt]
$+$ \texttt{Equations} &
quotients of free algebras &
finitely presentable universal algebra \\[3pt]
$+$ \texttt{Rewrites} &
theories with dynamics &
domain-specific languages \\[3pt]
$+$ parameters, and the algebra of theories &
theories built from other theories &
modules, in the sense of ML functors \\
\bottomrule
\end{tabular}
\end{center}

An ordinary data definition language stops on the first rung. Records, sums,
products, lists, maps: that is the whole of what most type systems can say about
data, and it is a lot. The second and third rungs are where this differs, and
the second rung is the one that will surprise a working developer most, so take
it first. The fourth is the subject of \S\ref{sec:reuse}--\S\ref{sec:module},
and answers the question every developer asks on seeing the declaration above:
\emph{can I reuse part of one of these?}

\subsection{Rung two: equations, or the comment your type system cannot read}

Every codebase of any size contains comments of this form:

\begin{quote}\small
\emph{This is a list, but order doesn't matter --- always compare as a set.}\\
\emph{This is a string, but compare case-insensitively.}\\
\emph{Addresses with an empty line-2 are equal to the same address without it.}
\end{quote}

Each is an equation. Each is being maintained by hand, in prose, because the
type system cannot express it. The consequences are familiar: a hand-written
\texttt{equals}, a normalisation function called on the way in and forgotten in
one code path, a cache keyed by the wrong representative, and a class of bug
that only appears when two spellings of the same value meet.

An \texttt{equations} block moves that comment into the definition:

\begin{lstlisting}[language=mettail]
Equations {
    (NQuote (PDrop N)) == N;
}
\end{lstlisting}

The payoff is not that the compiler now checks something. It is that
\emph{pattern matching happens up to the equations}. When the comprehension of
\S\ref{sec:cfl} matches \texttt{@spec} against a datum, it matches modulo the
theory that types the channel. A developer never writes normalisation, never
chooses a canonical form, and never has the bug where one code path forgot.

This is why the \texttt{PPar} constructor above takes a \texttt{HashBag} rather
than a list: parallel composition is associative and commutative, and rather
than being a special case hard-wired into an interpreter, that is a property of
the data declaration. The same mechanism is available for an invoice.

\begin{aside}[The reach of the second rung]
``Finitely presentable universal algebra'' is the mathematician's name for the
second rung, and it is worth translating: it means anything you can specify by
saying what the constructors are and which combinations of them are equal.
Monoids, groups, lattices, sets-as-quotients-of-lists, normalised paths,
case-insensitive strings, and commutative parallel composition are all on this
rung. Most of a data model's invisible invariants live here.
\end{aside}

\subsection{Rung three: rewrites, or data that wiggles}

The third rung adds \texttt{rewrites}, and with it the thing an ordinary schema
cannot do: the data moves on its own.

\begin{lstlisting}[language=mettail]
Rewrites {
    Beta     : (App (Lam ^x.body) arg) ~> (subst ^x.body arg);
    AppCongL : if M0 ~> M1 then (App M0 N) ~> (App M1 N);
}
\end{lstlisting}

That is the lambda calculus, in two of the four lines a real definition takes.
The first line is the computation rule --- and note that beta reduction is
literally the substitution of an argument into an abstraction, written as such.
The second is a congruence rule saying where the first is allowed to fire; the
\texttt{if \dots then} premise is the horizontal bar of an inference rule. Together they make a term of type
\texttt{Term} not an inert value but a small running program.

Once a channel is typed by such a definition, sending on that channel means
sending a term of a language --- a term that can be inspected, matched against,
and stepped by the recipient. Two agents that share a language definition can
exchange programs in a domain-specific language constructed for the purpose, and
the recipient knows what it received because the definition came with it.

The slogan for this rung is \emph{the data can wiggle}, and its practical form
is that the boundary between data and code, which most systems maintain by
convention and defend with a serialisation format, is a dial here rather than a
wall.

\subsection{Rung four: a language definition takes parameters}
\label{sec:reuse}

A developer who has just written the \texttt{RhoCalc} declaration above will
notice something about it. Parallel composition is associative, it is
commutative, and \texttt{PZero} is its unit. That is a commutative monoid. It is
also, verbatim, the structure of integer addition, of string concatenation up to
reordering, of set union, and of every other commutative monoid the developer
has ever implemented. Writing it out again is the same waste it always was.

The fourth rung is the answer, and it is the familiar one.

\begin{lstlisting}[language=mettail]
Theory EmptySet() {
    Types { Elem; }
    Exports { Elem; }
}

Theory Monoid(s: EmptySet) {
    s
      Terms {
        One  . |- "1" : Elem;
        Mult . x:Elem, y:Elem |- "(" x "*" y ")" : Elem;
      }
      Equations {
        (Mult (Mult x y) z) == (Mult x (Mult y z));
        (Mult x (One)) == x;
        (Mult (One) x) == x;
      }
}

Theory CommutativeMonoid(m: Monoid) {
    m
      Replacements {
        One  => Zero . |- "0" : Elem;
        Mult => Plus . x:Elem, y:Elem |- "(" x "+" y ")" : Elem;
      }
      Equations {
        (Plus x y) == (Plus y x);
      }
}
\end{lstlisting}

Two things are happening, and the second is the unfamiliar one.

\texttt{Monoid} takes a theory as a parameter and extends it. That is an ML
functor, or a generic class, or whatever the reader's language calls a module
parameterised by another module. Note where the parameter appears: not only in
the head, but as the \emph{first thing in the body}. \texttt{s} is the theory the
chain starts from, and each block after it applies to the result of the one
before. A body is not a record of five fields; it is a pipeline.

And \texttt{CommutativeMonoid} does something a generic class cannot: a
\texttt{Replacements} block \emph{renames an inherited constructor and gives it
new concrete syntax}. \texttt{Zero} and \texttt{Plus} are not new operations
sitting alongside \texttt{One} and \texttt{Mult}; they \emph{are} \texttt{One}
and \texttt{Mult}, spelled additively. Every equation proved of the one holds of
the other, because they are the same element of the same theory. This is the
distinction that the rest of \S\ref{sec:sharing} turns on, so it is worth
naming now: a replacement changes a spelling and preserves an identity.

Which is what lets \flang{}'s parallel composition be defined rather than
declared:

\begin{lstlisting}[language=mettail]
Theory ParMonoid(cm: u.CommutativeMonoid) {
    cm
      Exports {
        Elem => Proc;
      }
      Replacements {
        Zero => PZero . |- "0" : Proc;
        Plus => PPar  . ps:HashBag(Proc) |- "{" ps.*sep("|") "}" : Proc;
      }
      Rewrites {
        RPar : if S ~> T then (PPar {S, ...rest}) ~> (PPar {T, ...rest});
      }
}
\end{lstlisting}

Parallel composition is commutative-monoid multiplication with the carrier
renamed to \texttt{Proc}. That is the whole definition. Associativity,
commutativity, and the unit law arrive with it and are not restated.

\begin{aside}[Why the bag, again]
Notice what \texttt{PPar} did to arity on its way through. \texttt{Plus} is
binary. \texttt{PPar} takes a \texttt{HashBag}. A bag is precisely a sequence
quotiented by associativity and commutativity, so moving to a bag turns two of
the inherited equations from things checked during matching into things true of
the representation. This is the second rung and the first rung meeting: an
equation, made structural. It is also the one place where the surface currently
outruns its own checking, and \S\ref{sec:status} says so.
\end{aside}

\subsection{The module declaration}
\label{sec:module}

Theories do not float free. They are declared inside a \texttt{Module}, which is
the unit of distribution --- the thing that has a name, gets published to an
address, and is imported by other people's code.

\begin{lstlisting}[language=mettail]
import "UnivAlg.module" as u
import Monoid from "UnivAlg.module"

Module Rholang {

    Theory ParMonoid(cm: u.CommutativeMonoid) { ... }
    Theory QuoteDropCalc(pm: ParMonoid)       { ... }
    Theory RhoCalc(qd: QuoteDropCalc)         { ... }
    Theory NewReplCalc(qd: QuoteDropCalc)     { ... }

    Theory Rholang(nr: NewReplCalc, rc: RhoCalc) {
        nr \/ rc
    }

    Theory FreeRholang() {
        let s  = u.EmptySet() in (
        let m  = u.Monoid(s) in (
        let cm = u.CommutativeMonoid(m) in (
        let pm = ParMonoid(cm) in (
        let qd = QuoteDropCalc(pm) in (
        let nr = NewReplCalc(qd) in (
        let rc = RhoCalc(qd) in (
        Rholang(nr, rc)
        )))))))
    }

    theory FreeRholang()
}
\end{lstlisting}

The theory names are the historical ones and are worth reading as such:
\texttt{Rholang} here names the calculus \flang{} inherited its control flow
from, presented as one theory among others. Naming it is the point ---
\S\ref{sec:ddl} began by saying the calculus underneath is not privileged, and
this is what that looks like in a file.

Four things in that skeleton are worth calling out, because each is a decision a
developer will meet on the first afternoon.

\paragraph{Imports name an address, not a path.} \texttt{import "\dots" as u}
brings a whole module in under an alias, after which its theories are reached by
dotted path --- \texttt{u.CommutativeMonoid}. \texttt{import Monoid from "\dots"}
brings one theory in unqualified. The string is a URL. This is the choice that
aligns the module system with the deployment target: a language definition
published on chain has an address, and importing it is the same act whether it
came from the next directory or from another party's node. Resolution happens at
compile time, in the client toolchain. A node never fetches anything.

\paragraph{Capital \texttt{Theory} declares; lowercase \texttt{theory}
instantiates.} A module full of \texttt{Theory} declarations has produced no
language yet; it has produced a vocabulary for building one. The final
\texttt{theory FreeRholang()} is the entry point, and it is what the elaborator
runs. The last such line in the entry module wins. Keeping the two acts apart is
what allows a library module to ship a dozen theories and commit to none of
them.

\paragraph{A theory body is an expression, and it has a grammar.} Everything
between the braces of a \texttt{Theory} declaration is one expression built from
three kinds of thing.

\begin{center}
\small
\begin{tabular}{@{}L{0.20\textwidth}L{0.34\textwidth}L{0.36\textwidth}@{}}
\toprule
\textbf{Kind} & \textbf{Forms} & \textbf{What it does} \\
\midrule
atoms &
\texttt{Empty}, \texttt{free(Path)}, a parameter name, an application
\texttt{QuoteDropCalc(pm)}, \texttt{let x = e in (e)} &
starts a chain, or names an instance so it can be used twice \\[3pt]
postfix builders &
\texttt{Types}, \texttt{Exports}, \texttt{Replacements}, \texttt{Terms},
\texttt{Equations}, \texttt{Rewrites} &
each takes the theory to its left and returns a larger one; they chain, and
they may repeat \\[3pt]
combinators &
$/\!\backslash$ then $\backslash\!/$ then $\backslash$, tightest to loosest &
common fragment, combination, subtraction \\
\bottomrule
\end{tabular}
\end{center}

Because the builders form an ordered chain rather than a record, they can
interleave --- \texttt{AbelianGroup} in the universal-algebra module does a
\texttt{Replacements} block and then joins with a sibling --- and they carry an
obligation the record form did not: \emph{a \texttt{Terms} block must precede
any \texttt{Equations} or \texttt{Rewrites} block that mentions its labels}.
A forward reference is a diagnostic, not a resolution pass.

\paragraph{\texttt{Types} declares, \texttt{Exports} publishes.} These are two
jobs and they used to be one, which is worth a sentence because the separation
is what turns a class of modelling error into a compiler message.
\texttt{Types} brings a syntactic category into being. \texttt{Exports} controls
which of them a consumer can see, and optionally renames them on the way out
--- \texttt{Elem => Proc} above. A theory can now have a category that is real
and private. More to the point, two theories that each declare a category
called \texttt{Name} are now visibly declaring two different categories, which
is exactly the mistake \S\ref{sec:sharing} is about.

\begin{aside}[What a replacement does not need any more]
An earlier form of \texttt{Replacements} carried an integer profile ---
\texttt{[0,1]} --- to say how the replacement's arguments corresponded to the
original's. Term rules now name their arguments, so the correspondence is read
off the names and the profile is gone. Compare \texttt{Mult}'s context
\texttt{x:Elem, y:Elem} with \texttt{Plus}'s: identical names, identity
mapping. Naming them in the other order permutes them. This is a small thing
that removes an entire category of silent error.
\end{aside}

\subsection{What sharing means, and why the compiler argues about it}
\label{sec:sharing}

The combination operator has a subtlety that is worth a developer's attention,
because the compiler will raise it whether or not the developer was expecting
the conversation.

Look again at \texttt{FreeRholang}. One theory adds restriction and replication,
another adds send and receive, and both are built on \texttt{QuoteDropCalc}.
Combining them must not produce two parallel compositions, two \texttt{Name}
categories, and two quotation operators.

The \texttt{let} chain is what prevents it. \texttt{QuoteDropCalc(pm)} is
elaborated once and bound to \texttt{qd}; both \texttt{nr} and \texttt{rc} are
built from that one instance; so when they are combined, everything they
inherited from it is recognised as the same thing arriving by two routes, and
appears once. The combination is a pushout over the shared ancestor rather than
a union.

What makes that work mechanically is that every element gets an identity token
when it is introduced, and carries it through inheritance --- including through
a \texttt{Replacements} block, which is why renaming preserves sharing. Two
elements are the same element if they have the same origin. Combination
identifies elements by origin, the common fragment intersects by origin, and
subtraction removes by origin. Names have nothing to do with it.

The practical consequence is that \emph{sharing is a claim, and the compiler
holds you to it}. Build the two halves from two separate instances and combine
them, and the result is a diagnostic rather than a silently doubled calculus.
The reference elaborator ships a corpus of these; the four that a working
developer is most likely to trip over are worth having seen once:

\begin{center}
\small
\begin{tabular}{@{}L{0.30\textwidth}L{0.62\textwidth}@{}}
\toprule
\textbf{Diagnostic} & \textbf{What you did, and what it means} \\
\midrule
\texttt{[join-collision]} &
combined two theories that each introduced a label or a category
independently. They are two things with one spelling. Either they should have
descended from a shared instance, or they should have different names \\[3pt]
\texttt{[unshared-name]} &
the same, in its most common disguise: two independently instantiated copies of
the same parameter theory passed into one combination \\[3pt]
\texttt{[repeat-label]} &
declared a label twice in one theory. A replacement retargets an existing
label; it does not introduce a second one \\[3pt]
\texttt{[forward-reference]} &
an \texttt{Equations} or \texttt{Rewrites} block mentioned a label its own
chain had not yet introduced \\
\bottomrule
\end{tabular}
\end{center}

This is unfamiliar and it is worth sitting with, because it is the same
diagnostic that catches genuine modelling errors. A rig has one carrier and two
monoid structures on it --- addition and multiplication --- and the natural way
to write it is to instantiate \texttt{Monoid} once and pass it as both
arguments. That is wrong: it asserts that addition and multiplication are the
same operation. The compiler says so. Instantiating \texttt{Monoid} twice over
one shared \texttt{EmptySet} is the correct structure, and it says exactly what
a rig is: two monoids, one carrier.

A developer used to inheritance will recognise the diamond problem here, and
should notice that the resolution is different in kind. There is no linearised
method resolution order and no ``last one wins''. Two elements are the same
element if they have the same origin, and otherwise they are two elements; if
they also have the same name, that is an error rather than a silent override.

\subsection{The parameter as a contract: extending a theory you did not write}
\label{sec:loadtest}

The reuse story so far is about not repeating yourself. There is a second
payoff, and for anyone maintaining a language definition alongside other people
it is the larger one.

Consider a design that has an object language --- a dependently typed calculus
with typed holes, used as the specification format for a marketplace of partly
specified programs~\cite{PF} --- and a second design that adds prices to it. The
priced language is not a copy of the first with extra lines. It is a theory that
takes the first as a parameter:

\begin{lstlisting}[language=mettail]
Theory PricedPFLam(pf: PFLam) {
    pf
      Types {
        Price;       // an element of the preference quantale
        Obsv;        // a time-varying observable
      }
      Terms {
        Priced . h:HoleId, ty:Term, p:Price
               |- "?" h ":" ty "@" p                  : Term;
        Scale  . o:Obsv, c:Term |- "scale" "(" o "," c ")" : Term;
        Give   . c:Term         |- "give" "(" c ")"        : Term;
      }
      Equations {
        (Give (Give c)) == c;
      }
}
\end{lstlisting}

Everything established about \texttt{PFLam} arrives with the parameter. Nothing
is restated. And --- the point --- if a later revision of \texttt{PFLam} renames
\texttt{Hole} or changes its arity, the break is reported \emph{here}, at
elaboration, rather than discovered later by a reader noticing that two
documents disagree. Before the fourth rung existed, this delta was written as
prose on a listing: \emph{only additions shown}. A prose annotation cannot be
checked. A parameter can.

The same applies at combination. If a venue also imports a reputation theory
over the same hole identifiers, the two join over their shared \texttt{PFLam}
instance and the identifiers are identified once. Passing two independently
instantiated copies is rejected, and correctly: it would assert that the priced
holes and the reputed holes are different holes.

\subsection{Channels are typed by language definitions}

The connection back to \S\ref{sec:cfl} is one sentence: \emph{a channel's type
is a language definition} --- which is to say, an elaborated \texttt{Theory}.

That single design decision explains the rest of the note. Sending on the
channel means sending a term of that theory. Matching a pattern on it means
matching a term of that theory, up to that theory's equations. And --- the point
of \S\ref{sec:logic} --- asking a condition about what arrived means asking a
question in a logic that is \emph{generated from} that theory, because the
theory says what there is to ask about.

For the forager, this means the specimen language is a language definition,
\texttt{Chi} is a formula in the logic that definition generates, and there was
never a moment at which someone hand-designed a predicate vocabulary.

\subsection{MeTTa, and what this adds to it}
\label{sec:metta}

Readers coming from SingularityNET's Hyperon stack~\cite{Hyperon} will have
recognised the shape of the third rung, because MeTTa is built on it. The relationship is worth
stating precisely, in both directions.

\begin{center}
\small
\begin{tabular}{@{}L{0.24\textwidth}L{0.30\textwidth}L{0.38\textwidth}@{}}
\toprule
\textbf{MeTTa} & \textbf{\flang{}} & \textbf{Relationship} \\
\midrule
rules &
the \texttt{Rewrites} block &
subsumed; MeTTa's rules are rewrites in this sense \\[3pt]
atom spaces &
channels &
subsumed, and typed: a channel carries terms of a declared language, an
atom space carries atoms \\[3pt]
--- &
the \texttt{Equations} block &
\textbf{absent in MeTTa}: no equational layer, so every quotient must be
hand-coded as rules, and confluence becomes the programmer's problem \\[3pt]
query, then update &
one witnessed communication &
\textbf{absent in MeTTa}: query and update are separate acts, so there is no
atomic multi-party meeting and no witness \\
\bottomrule
\end{tabular}
\end{center}

The two gaps compound. Without equations, a developer who wants ``these two
atoms are the same'' writes rewrite rules that say so, and then owns the
confluence question --- whether the order in which rules fire can change the
answer. Without a witnessed transaction, a developer who wants two atom spaces
updated together writes a protocol, and then owns the partial-failure question.
Both are exactly the questions the two blocks and the communication rule
answer for free.

None of which is a dismissal. MeTTa has a large working ecosystem, a substantial
body of practice around the atom space as a shared knowledge medium, and
Probabilistic Logic Networks~\cite{PLN}, which is a serious piece of engineering for
reasoning under uncertainty. It is worth saying plainly that PLN is not left
behind here: in \S\ref{sec:graded} it comes back as one of the value algebras a
\texttt{where} clause can be graded in, which is to say that PLN inference can
drive execution rather than sitting beside it. The relationship this note is
proposing is embedding with an addition, not replacement.

\subsection{What an operation costs}
\label{sec:cost}

One more piece belongs to the data definition story, because it is the thing
that makes the forager's budget real rather than notional.

Cost accounting in this system adds exactly two surface forms, and both of them
lower to ordinary \flang{}:

\begin{center}
\small
\begin{tabular}{@{}lll@{}}
\toprule
\textbf{Surface} & \textbf{Lowers to} & \textbf{Which is} \\
\midrule
\texttt{s :: ()} & \texttt{$\Sigma\sem{s}$!(Nil)} & a send: mint a token \\
\texttt{\{\% P \%\}[s]} & \texttt{for(t <- $\Sigma\sem{s}$)\{*t | P\}} & a receipt: run \texttt{P} once a token arrives \\
\bottomrule
\end{tabular}
\end{center}

That is the entire mechanic. Spending fuel \emph{is} a communication: the gate
consumes one token, releases its payload, and then \texttt{P} runs. If no token
is available on the signature channel, the gate parks forever and \texttt{P}
never runs --- and that is what metering means here. A budget is a stack of
nested sends, and its depth is how many operations it funds.

Two things follow that a developer should notice.

First, metering is not a virtual-machine feature bolted onto a language. It is a
program \emph{in} the language, and the transpiler's job is only to write it for
you. Anything true of unmetered \flang{} is true of metered \flang{}.

Second, it composes with Claim~\ref{claim:witness} in a satisfying way. A
communication that is not funded is a communication that does not happen, and it
is indistinguishable from a communication whose guard failed. Both are
non-events. There is no separate failure mode for running out of budget, no
exception, and nothing to unwind --- which is exactly the property that lets the
forager's death be modelled as ``it stops'' rather than as an error path.

The repository~\cite{f1r3node} ships a worked demonstration of this at scale: a supply chain in
which factories produce, carriers ship, sellers wholesale and buyers retail,
with money and inventory both conserved, and every single operation gated by a
token. It is the best available answer to ``show me this doing something real''.

\subsection{What actually runs}
\label{sec:status}

Honesty about status is more useful to a developer than enthusiasm, so the
table below separates three things a reader has every right to keep apart:
\emph{specified} means the surface is frozen and written down; \emph{implemented}
means there is code a reader can run; \emph{planned} means neither yet.

\begin{center}
\small
\begin{tabular}{@{}L{0.40\textwidth}L{0.52\textwidth}@{}}
\toprule
\textbf{Capability} & \textbf{Status} \\
\midrule
\texttt{Theory} declarations: types, terms, equations, rewrites; generated
parser, printer, substitution, rewrite engine &
implemented, in the reference elaborator and in
\texttt{mettail-rust}~\cite{mettailrust}, with language definitions shipped for
lambda, ambient, the \rhoc{} calculus and a calculator \\[3pt]
The fourth rung (\S\ref{sec:reuse}, \S\ref{sec:module}): theory parameters,
\texttt{Types}, \texttt{Exports}, \texttt{Replacements}, the algebra of
theories, and pushout sharing (\S\ref{sec:sharing}) &
implemented, as the \texttt{mettail-elab} crate on the
\texttt{feature/module-syntax} branch of
\texttt{f1r3node-rust}~\cite{modulesyntax}. Dependency-free, with the
universal-algebra tower and the \rhoc{} calculus as its worked corpus and a
corpus of negative cases for the diagnostics of \S\ref{sec:sharing} \\[3pt]
The same surface in the \emph{live} parser &
\emph{not yet}. The grammar delta for the tree-sitter parser sits beside the
elaborator as a patch rather than an applied change, and the elaborator is not
yet on the node's workspace build. This is the next step, and it is the honest
boundary of this section \\[3pt]
\texttt{Module}, and imports by URL &
specified and implemented in the elaborator; resolution is a compile-time step
in the client toolchain, so a node never performs it \\[3pt]
\texttt{where} guards on receipts, and \texttt{where} guards on \texttt{match}
cases with fall-through &
implemented on the \texttt{cost-accounting-transpiler} branch of
\texttt{f1r3node-rust}~\cite{f1r3node}, with runnable examples and compile tests \\[3pt]
Cost accounting: minting, gates, compound signatures, budget stacks &
implemented on the same branch, with a transpiler and a documented demo \\[3pt]
Arity-changing replacements (the bag in \texttt{ParMonoid}) &
accepted by the elaborator, \emph{not} checked. Replacing a binary constructor
by a collection-valued one leaves the inherited binary equations written over a
constructor that no longer has that shape. Either replacements should be
arity-preserving, or there should be a way to say that an inherited equation has
become structural. Open \\[3pt]
Conditions drawn from the \emph{generated logic} (\S\ref{sec:logic}) rather than
from ordinary boolean expressions &
planned. The guard slot exists and is evaluated before commit, but a condition
today is an ordinary boolean expression; the structural and behavioural
connectives are specified in the research notes, and connecting an elaborated
theory to the guard is what would close the gap \\[3pt]
Graded conditions and simulation (\S\ref{sec:graded}) &
implemented as a separate simulator crate~\cite{WeightedGSLT}, with the
stochastic and quantum readings both built and tested; not part of the node's execution path, by
design (\S\ref{sec:simnotinterp}) \\
\bottomrule
\end{tabular}
\end{center}

The condition language a developer can write \emph{today} is boolean expressions
over the variables the patterns bound --- \texttt{x > 0}, and so on. Everything
in the next section is about what that slot is going to hold, and why the choice
is more interesting than it looks.
